\documentclass[11pt]{article}
% ===========================================================================
%  Shared preamble for all MPM2D worksheets — input right after \documentclass.
%  (Files beginning with "_" are skipped by mpm2d-worksheet-publish.mjs.)
% ===========================================================================
\usepackage[letterpaper,top=2.2cm,bottom=1.9cm,left=1.6cm,right=1.6cm,headheight=28pt,footskip=24pt]{geometry}
\usepackage{amsmath,amssymb}
\usepackage[T1]{fontenc}
\usepackage[utf8]{inputenc}
\usepackage{lmodern}
\usepackage{xcolor}
\usepackage[most]{tcolorbox}
\usepackage{enumitem}
\usepackage{multicol}
\usepackage{fancyhdr}
\usepackage{tikz}
\usetikzlibrary{calc}
\usepackage{pgfplots}
\pgfplotsset{compat=1.18}

\definecolor{exblue}{HTML}{4A90E2}
\definecolor{exbg}{HTML}{E6F3FF}
\definecolor{qorange}{HTML}{E69138}
\definecolor{qbg}{HTML}{FFF7CC}
\definecolor{keygray}{HTML}{475569}
\definecolor{keybg}{HTML}{F1F5F9}
\definecolor{iagreen}{HTML}{14653B}
\definecolor{titleblue}{HTML}{1D4ED8}
% Rotating palette so each worked example gets its own colour.
\definecolor{ex0}{HTML}{2563A0}\definecolor{ex1}{HTML}{3B7D3B}\definecolor{ex2}{HTML}{8A5A00}
\definecolor{ex3}{HTML}{A3327A}\definecolor{ex4}{HTML}{6D28A3}\definecolor{ex5}{HTML}{B08900}
\definecolor{ex6}{HTML}{0E7490}\definecolor{ex7}{HTML}{9A3412}\definecolor{ex8}{HTML}{15803D}

\pagestyle{fancy}
\fancyhf{}
\fancyhead[L]{\footnotesize NAME: \underline{\hspace{3.4cm}}}
\fancyhead[C]{\textbf{Integration Academy}\\[-2pt]{\scriptsize Math Simplified \textperiodcentered{} Success Amplified}}
\fancyhead[R]{\footnotesize MPM2D}
\fancyfoot[L]{\scriptsize dr.merie@IntegrationAcademy.ca}
\fancyfoot[C]{\scriptsize\thepage}
\fancyfoot[R]{\scriptsize IntegrationAcademy.ca}
\renewcommand{\headrulewidth}{1.2pt}
\renewcommand{\footrulewidth}{0.4pt}
\setlength{\parindent}{0pt}
\setlength{\parskip}{4pt}

% Examples are NOT breakable, so each example + its solution stays on one page.
\newtcolorbox{example}[2]{colframe=#1,colback=#1!7,coltitle=white,colbacktitle=#1,fonttitle=\bfseries,title={#2},arc=3pt,boxrule=1pt}
\newtcolorbox{qbox}[1]{colframe=qorange,colback=qbg,coltitle=white,colbacktitle=qorange,fonttitle=\bfseries,title={#1},arc=3pt,breakable,boxrule=1pt}
\newtcolorbox{keybox}{colframe=keygray,colback=keybg,coltitle=white,colbacktitle=keygray,fonttitle=\bfseries,title={$\checkmark$\ Answer Key},arc=3pt,breakable,boxrule=1pt}
\newtcolorbox{recapbox}{colframe=keygray,colback=keybg,coltitle=white,colbacktitle=keygray,fonttitle=\bfseries,title={Question Recap},arc=3pt,breakable,boxrule=1pt}
\newtcolorbox{ideabox}{colframe=titleblue,colback=titleblue!6,coltitle=white,colbacktitle=titleblue,fonttitle=\bfseries,title={The Big Ideas},arc=3pt,breakable,boxrule=1.2pt}
\newtcolorbox{learnbox}{colframe=iagreen,colback=iagreen!5,coltitle=white,colbacktitle=iagreen,fonttitle=\bfseries,title={Core Concepts},arc=3pt,breakable,boxrule=1.2pt}
\newcommand{\lh}[1]{\par\smallskip{\bfseries\color{iagreen}#1}\par\nobreak\smallskip}

\newcommand{\soln}{\par\smallskip\textit{\textbf{Solution.}}\ }
\newcommand{\workspace}[1]{\par\vspace{3pt}\fcolorbox{black!30}{white}{\begin{minipage}[t][#1][t]{\dimexpr\linewidth-2\fboxsep\relax}\ \end{minipage}}\par}
\newcommand{\sech}[1]{\par\vspace{8pt}{\Large\bfseries\color{iagreen}#1}\par\nobreak\vspace{2pt}{\color{black!20}\hrule height1pt}\vspace{6pt}}

% A compact coordinate plot sized for an example box: \eplot{xmin}{xmax}{ymin}{ymax}{body}
\newcommand{\eplot}[5]{\begin{center}\begin{tikzpicture}
\begin{axis}[width=6.8cm,height=4.4cm,axis lines=middle,xlabel={\tiny$x$},ylabel={\tiny$y$},
grid=both,grid style={gray!16},xmin=#1,xmax=#2,ymin=#3,ymax=#4,
xtick distance=2,ytick distance=2,tick label style={font=\tiny},axis line style={-},samples=120]
#5
\end{axis}\end{tikzpicture}\end{center}}

% Same as \eplot but with its own tick spacing: \eplotx{xmin}{xmax}{ymin}{ymax}{xtick}{ytick}{body}
\newcommand{\eplotx}[7]{\begin{center}\begin{tikzpicture}
\begin{axis}[width=8.4cm,height=5.4cm,axis lines=middle,xlabel={\tiny$x$},ylabel={\tiny$y$},
grid=both,grid style={gray!16},xmin=#1,xmax=#2,ymin=#3,ymax=#4,
xtick distance=#5,ytick distance=#6,tick label style={font=\tiny},axis line style={-},samples=120]
#7
\end{axis}\end{tikzpicture}\end{center}}

% A sine-curve plot (degrees on x, 0..360): \splot{ymin}{ymax}{body}
\newcommand{\splot}[3]{\begin{center}\begin{tikzpicture}
\begin{axis}[width=8.6cm,height=4.6cm,axis lines=middle,xlabel={\tiny$x\,(^\circ)$},ylabel={\tiny$y$},
grid=both,grid style={gray!16},xmin=0,xmax=360,ymin=#1,ymax=#2,
xtick distance=90,ytick distance=1,tick label style={font=\tiny},axis line style={-},samples=180]
#3
\end{axis}\end{tikzpicture}\end{center}}

% A small coordinate plot: \plot{xmin}{xmax}{ymin}{ymax}{body}
\newcommand{\plot}[5]{\begin{center}\begin{tikzpicture}
\begin{axis}[width=8.4cm,height=6cm,axis lines=middle,xlabel={\scriptsize$x$},ylabel={\scriptsize$y$},
grid=both,grid style={gray!16},xmin=#1,xmax=#2,ymin=#3,ymax=#4,
xtick distance=2,ytick distance=2,tick label style={font=\tiny},axis line style={-}]
#5
\end{axis}\end{tikzpicture}\end{center}}

% Right triangle (angle theta at bottom-left, right angle at bottom-right):
%   \rtri{drawBase}{drawHeight}{oppLabel}{adjLabel}{hypLabel}
\newcommand{\rtri}[5]{\begin{center}\begin{tikzpicture}[scale=0.85]
\coordinate (A) at (0,0);\coordinate (B) at (#1,0);\coordinate (C) at (#1,#2);
\draw[thick,exblue] (A)--(B)--(C)--cycle;
\draw[black] ($(B)+(-0.32,0)$)--($(B)+(-0.32,0.32)$)--($(B)+(0,0.32)$);
\node[below] at ($(A)!0.5!(B)$){#4};
\node[right] at ($(B)!0.5!(C)$){#3};
\node[above left] at ($(A)!0.5!(C)$){#5};
\node at (0.7,0.22){$\theta$};
\end{tikzpicture}\end{center}}

% General (oblique) triangle with vertex labels and opposite-side labels:
%   \gtri{Alabel}{Blabel}{Clabel}{aSide}{bSide}{cSide}
\newcommand{\gtri}[6]{\begin{center}\begin{tikzpicture}[scale=0.85]
\coordinate (A) at (0,0);\coordinate (B) at (5,0);\coordinate (C) at (1.6,3);
\draw[thick,exblue] (A)--(B)--(C)--cycle;
\node[below left] at (A){#1};\node[below right] at (B){#2};\node[above] at (C){#3};
\node[below] at ($(A)!0.5!(B)$){#6};
\node[right] at ($(B)!0.5!(C)$){#4};
\node[left] at ($(A)!0.5!(C)$){#5};
\end{tikzpicture}\end{center}}

\fancyhead[R]{\footnotesize MHF4U \textperiodcentered{} 4.4}
\begin{document}

\begin{center}
{\LARGE\bfseries\color{titleblue}Worksheet 4.4 --- Applications of Exponential \& Log Models}\\[2pt]
{\small Grade 12 Advanced Functions (MHF4U) \textperiodcentered{} Unit 4: Exponential and Logarithmic Functions}
\end{center}

\begin{ideabox}
Model with $A=A_0\,b^{t/p}$ (growth $b>1$, decay $0<b<1$); use logs for doubling time and half-life; log scales (pH, dB, Richter) step by factors of ten.
\begin{itemize}[leftmargin=*,itemsep=2pt]
  \item $A=A_0\,b^{t/p}$, where $p$ is the doubling time or half-life.
  \item Doubling time $=\dfrac{\log2}{\log(\text{growth factor})}$.
  \item $\text{pH}=-\log[\text{H}^+]$; each Richter step is $\times10$.
\end{itemize}
\end{ideabox}

\begin{learnbox}
\lh{Growth and decay}Exponential models $A=A_0 b^{\,t}$ describe populations, radioactive decay, and cooling; logarithms solve them for the time.
\lh{Compound interest}Investments grow by $A=P(1+i)^n$, and finding the number of periods $n$ requires a logarithm.
\lh{Logarithmic scales}Quantities spanning huge ranges — sound in decibels, earthquakes on the Richter scale, acidity as pH — use \textbf{logarithmic} scales.
\end{learnbox}

\sech{Worked Examples}

\begin{example}{ex0}{Example 1 --- Growth}
$P=120\cdot2^{t/2}$. Find $P$ at $t=6$.\soln $P=120\cdot2^{3}=960$. Growth $2^x$ vs decay $0.5^x$:\eplot{-3}{3}{-1}{9}{\addplot[exblue,very thick,domain=-3:3.1,samples=80]{2^x};\addplot[qorange,very thick,domain=-3.1:3,samples=80]{0.5^x};}
\end{example}

\begin{example}{ex1}{Example 2 --- Half-life}
$A=96(\tfrac12)^{t/2}$. Find $A$ at $t=6$.\soln $A=96\cdot(\tfrac12)^3=96\cdot\tfrac18=12$ mg.
\end{example}

\begin{example}{ex2}{Example 3 --- Doubling time}
How long to double at 3\% /yr?\soln $2=1.03^t\Rightarrow t=\dfrac{\log2}{\log1.03}\approx23.4$ yr.
\end{example}

\begin{example}{ex3}{Example 4 --- pH scale}
Find the pH if $[\text{H}^+]=10^{-6}$.\soln $\text{pH}=-\log(10^{-6})=6$.
\end{example}

\begin{example}{ex4}{Example 5 --- Richter scale}
How much stronger is magnitude 8 than magnitude 7?\soln Difference $1\Rightarrow10^1=10\times$.
\end{example}

\begin{example}{ex5}{Example 6 --- Growth}
$P=250\cdot2^{t/4}$. Find $P$ at $t=12$.\soln $250\cdot2^3=2000$.
\end{example}

\begin{example}{ex6}{Example 7 --- Half-life}
$A=200(\tfrac12)^{t/6}$. Find $A$ at $t=12$.\soln $200\cdot\tfrac14=50$.
\end{example}

\begin{example}{ex7}{Example 8 --- Doubling time}
How long to double at 6\% /yr?\soln $\dfrac{\log2}{\log1.06}\approx11.9$ yr.
\end{example}

\begin{example}{ex8}{Example 9 --- pH scale}
Find the pH if $[\text{H}^+]=10^{-11}$.\soln pH $=11$.
\end{example}

\clearpage
\sech{Your Turn --- Show Your Work}
For each question, show all steps clearly in the space provided.

\begin{qbox}{Question 1}
$P=300\cdot2^{t/6}$. Find $P$ at $t=12$.\workspace{2.4cm}
\end{qbox}

\begin{qbox}{Question 2}
$A=240(\tfrac12)^{t/4}$. Find $A$ at $t=8$.\workspace{2.4cm}
\end{qbox}

\begin{qbox}{Question 3}
How long to double at 8\% /yr (2 d.p.)?\workspace{2.4cm}
\end{qbox}

\begin{qbox}{Question 4}
pH if $[\text{H}^+]=10^{-3}$?\workspace{2.4cm}
\end{qbox}

\begin{qbox}{Question 5}
How much stronger is magnitude 6 than magnitude 3?\workspace{2.4cm}
\end{qbox}

\begin{qbox}{Question 6}
$P=100\cdot2^{t/10}$. Find $P$ at $t=20$.\workspace{2.4cm}
\end{qbox}

\begin{qbox}{Question 7}
$A=64(\tfrac12)^{t/2}$. Find $A$ at $t=6$.\workspace{2.4cm}
\end{qbox}

\begin{qbox}{Question 8}
How long to triple at 5\% /yr (2 d.p.)?\workspace{2.4cm}
\end{qbox}

\begin{qbox}{Question 9}
pH if $[\text{H}^+]=10^{-9}$?\workspace{2.4cm}
\end{qbox}

\begin{qbox}{Question 10}
A quake is $10^4$ times stronger than another. What is the magnitude difference?\workspace{2.4cm}
\end{qbox}

\begin{qbox}{Question 11}
$P=50\cdot3^{t/5}$. Find $P$ at $t=5$.\workspace{2.4cm}
\end{qbox}

\begin{qbox}{Question 12}
How long for 5\% /yr growth to reach 1.5 times the start (2 d.p.)?\workspace{2.4cm}
\end{qbox}

\begin{qbox}{Question 13 --- Challenge}
Bacteria double every 3 h from 500. Write the model and find the count at $t=9$ h.\workspace{3.5cm}
\end{qbox}

\clearpage
\sech{Question Recap \& Answer Key}

\begin{recapbox}
\begin{multicols}{2}
\small
\begin{enumerate}[leftmargin=*,itemsep=3pt]
  \item $P=300\cdot2^{t/6}$. Find $P$ at $t=12$.
  \item $A=240(\tfrac12)^{t/4}$. Find $A$ at $t=8$.
  \item How long to double at 8\% /yr (2 d.p.)?
  \item pH if $[\text{H}^+]=10^{-3}$?
  \item How much stronger is magnitude 6 than magnitude 3?
  \item $P=100\cdot2^{t/10}$. Find $P$ at $t=20$.
  \item $A=64(\tfrac12)^{t/2}$. Find $A$ at $t=6$.
  \item How long to triple at 5\% /yr (2 d.p.)?
  \item pH if $[\text{H}^+]=10^{-9}$?
  \item A quake is $10^4$ times stronger than another. What is the magnitude difference?
  \item $P=50\cdot3^{t/5}$. Find $P$ at $t=5$.
  \item How long for 5\% /yr growth to reach 1.5 times the start (2 d.p.)?
  \item Bacteria double every 3 h from 500. Write the model and find the count at $t=9$ h.
\end{enumerate}
\end{multicols}
\end{recapbox}

\begin{keybox}
\begin{multicols}{2}
\small
\begin{enumerate}[leftmargin=*,itemsep=3pt]
  \item $1200$
  \item $60$
  \item $\tfrac{\log2}{\log1.08}\approx9.01$ yr
  \item $3$
  \item $1000\times$
  \item $400$
  \item $8$
  \item $\tfrac{\log3}{\log1.05}\approx22.52$ yr
  \item $9$
  \item $4$
  \item $150$
  \item $\tfrac{\log1.5}{\log1.05}\approx8.31$ yr
  \item $P=500\cdot2^{t/3}$, $P(9)=4000$
\end{enumerate}
\end{multicols}
\end{keybox}

\end{document}
